% GENERATED FILE. Edit canonical lessons, then run npm run generate:books.
% canonical-source-hash: fnv1a64:b282341427ddcaa9
% canonical-lessons: TA-C287-poka-ventum, TA-C287-vantatal, TA-R287-advice-and-reasons

\chapter{Must, Must Not and Because}
\label{ch:ta-must-must-not-and-because}
\hlchaptermodality{287}

\begin{chapteropening}
\textbf{By the end of this chapter:} \emph{I can say what someone must, must not or need not do with pōka vēṇṭum, pōkak kūṭātu and vara vēṇṭām, and give a reason with vantatāl, ēṉeṉṟāl and kāraṇam.}

Complete the last lesson of chapter 287: give advice with must and must not, and the reason for it, before or after the result.
\end{chapteropening}

\section[\texorpdfstring{\ta{போக} \ta{வேண்டும்} (pōka vēṇṭum)}{போக வேண்டும் (pōka vēṇṭum)}]{\ta{போக} \ta{வேண்டும்} (pōka vēṇṭum) — must, must not, no need}

\label{lesson:TA-C287-poka-ventum}



\begin{quote}
Say \textquotedblleft{}I want water\textquotedblright{}. (\emph{eṉakkut taṇṇīr vēṇṭum}.) Then \textquotedblleft{}let's go\textquotedblright{}. (\emph{pōkalām}.)

Both hold a piece this lesson takes out.
\end{quote}

\subsection*{You'll want to know: the to-form and \ta{வேண்டும்}}
\textbf{\ta{போகலாம்}} is \textbf{\ta{போக}}, \textquotedblleft{}to go\textquotedblright{}, with \textbf{-\ta{லாம்}}, \textquotedblleft{}may\textquotedblright{}. \textbf{\ta{போக}} is the \textbf{to-form} of \textbf{\ta{போ}}. Put \textbf{\ta{வேண்டும்}} after it and you have \textquotedblleft{}must\textquotedblright{}:

\begin{quote}
\textbf{\ta{நான்} \ta{இன்று} \ta{மருத்துவரைப்} \ta{பார்க்க} \ta{வேண்டும்}.}
\end{quote}

\begin{quote}
\emph{nāṉ iṉṟu maruttuvaraip pārkka vēṇṭum.} — \textquotedblleft{}I have to see a doctor today.\textquotedblright{}
\end{quote}

The same \textbf{\ta{வேண்டும்}} as \textbf{\ta{எனக்குத்} \ta{தண்ணீர்} \ta{வேண்டும்}}: there \textquotedblleft{}is wanted\textquotedblright{} takes a thing, and the one who wants takes \textbf{-\ta{க்கு}}. Here it takes an action, and the person stands in front plainly: \textbf{\ta{நான்}}, \textbf{\ta{நீங்கள்}}.

Learn each to-form with its verb:

\begin{itemize}
\raggedright
  \item \textbf{\ta{போ}} $\to$ \textbf{\ta{போக}} (\emph{pōka}); \textbf{\ta{வா}} $\to$ \textbf{\ta{வர}} (\emph{vara}); \textbf{\ta{இரு}} $\to$ \textbf{\ta{இருக்க}} (\emph{irukka})
  \item \textbf{\ta{பார்}} $\to$ \textbf{\ta{பார்க்க}} (\emph{pārkka}); \textbf{\ta{படி}} $\to$ \textbf{\ta{படிக்க}} (\emph{paṭikka}); \textbf{\ta{கொடு}} $\to$ \textbf{\ta{கொடுக்க}} (\emph{koṭukka})
  \item \textbf{\ta{செய்}} $\to$ \textbf{\ta{செய்ய}} (\emph{ceyya}); \textbf{\ta{சாப்பிடு}} $\to$ \textbf{\ta{சாப்பிட}} (\emph{cāppiṭa}); \textbf{\ta{பேசு}} $\to$ \textbf{\ta{பேச}} (\emph{pēca})
\end{itemize}

In speech, \textbf{\ta{போக} \ta{வேண்டும்}} shrinks to \textbf{\ta{போகணும்}}; the written form keeps the long one.

\subsection*{You'll want to know: \ta{கூடாது} and \ta{வேண்டாம்}}
\begin{quote}
\textbf{\ta{இரவில்} \ta{தனியாகப்} \ta{போகக்} \ta{கூடாது}.}
\end{quote}

\begin{quote}
\emph{iravil taṉiyākap pōkak kūṭātu.} — \textquotedblleft{}You must not go alone at night.\textquotedblright{}
\end{quote}

\textbf{\ta{கூடாது}} (\emph{kūṭātu}, \textquotedblleft{}must not, it is not right\textquotedblright{}) forbids, and it does not change for the person. After the to-form the first letter doubles in writing: \textbf{\ta{போகக்} \ta{கூடாது}}, \textbf{\ta{சாப்பிடக்} \ta{கூடாது}}.

\begin{quote}
\textbf{\ta{நீங்கள்} \ta{நாளை} \ta{வர} \ta{வேண்டாம்}.}
\end{quote}

\begin{quote}
\emph{nīṅkaḷ nāḷai vara vēṇṭām.} — \textquotedblleft{}You don't need to come tomorrow.\textquotedblright{}
\end{quote}

\textbf{\ta{வேண்டாம்}}, \textquotedblleft{}is not wanted\textquotedblright{}, releases: \textquotedblleft{}no need to\textquotedblright{}, or a gentle \textquotedblleft{}don't\textquotedblright{}. \textbf{\ta{கூடாது}} is a ban; \textbf{\ta{வேண்டாம்}} is advice.

\subsection*{Your turn}
Say these aloud:

\begin{itemize}
\raggedright
  \item \textquotedblleft{}I have to see a doctor today\textquotedblright{} — \emph{nāṉ iṉṟu maruttuvaraip pārkka vēṇṭum}
  \item \textquotedblleft{}You must not go alone at night\textquotedblright{} — \emph{iravil taṉiyākap pōkak kūṭātu}
  \item \textquotedblleft{}You don't need to come tomorrow\textquotedblright{} — \emph{nīṅkaḷ nāḷai vara vēṇṭām}
  \item one thing you must do this week, and one thing you must not
\end{itemize}

\begin{tcolorbox}[breakable,colback=teal!4,colframe=teal!35!black,title={Before you move on}]
\textquotedblleft{}I must go\textquotedblright{}? (\textbf{\ta{நான்} \ta{போக} \ta{வேண்டும்}}.) \textquotedblleft{}You must not eat\textquotedblright{}? (\textbf{\ta{சாப்பிடக்} \ta{கூடாது}}.) \textquotedblleft{}No need to come\textquotedblright{}? (\textbf{\ta{வர} \ta{வேண்டாம்}}.) Which one forbids, and which one advises? (\textbf{\ta{கூடாது}} forbids; \textbf{\ta{வேண்டாம்}} advises.)
\end{tcolorbox}
\section[\texorpdfstring{\ta{வந்ததால்} (vantatāl)}{வந்ததால் (vantatāl)}]{\ta{வந்ததால்} (vantatāl) — because it happened}

\label{lesson:TA-C287-vantatal}



\begin{quote}
Say \textquotedblleft{}I have to see a doctor today.\textquotedblright{} (\emph{nāṉ iṉṟu maruttuvaraip pārkka vēṇṭum}.) Then \textquotedblleft{}So I stayed at home.\textquotedblright{} (\emph{ataṉāl nāṉ vīṭṭil iruntēṉ}.)
\end{quote}

\subsection*{You'll want to know: -\ta{ததால்}, 'because (it) did'}
\begin{quote}
\textbf{\ta{மழை} \ta{பெய்ததால்} \ta{நான்} \ta{வரவில்லை}.}
\end{quote}

\begin{quote}
\emph{maḻai peytatāl nāṉ varavillai.} — \textquotedblleft{}Because it rained, I did not come.\textquotedblright{}
\end{quote}

\textbf{\ta{பெய்ததால்}} is three pieces. \textbf{\ta{பெய்த}} is the past with \textbf{-\ta{அ}}, the piece you met in \textbf{\ta{வந்த} \ta{போது}}; \textbf{\ta{அது}}, \textquotedblleft{}it\textquotedblright{}, turns the event into a thing; \textbf{-\ta{ஆல்}}, \textquotedblleft{}because of\textquotedblright{}, is the \textbf{-\ta{ஆல்}} of \textbf{\ta{அதனால்}}. \textquotedblleft{}Because of its having rained.\textquotedblright{}

\begin{quote}
\textbf{\ta{எனக்குக்} \ta{காய்ச்சல்} \ta{இருந்ததால்} \ta{நான்} \ta{பள்ளிக்குப்} \ta{போகவில்லை}.}
\end{quote}

\begin{quote}
\emph{eṉakkuk kāyccal iruntatāl nāṉ paḷḷikkup pōkavillai.} — \textquotedblleft{}Because I had a fever, I did not go to school.\textquotedblright{}
\end{quote}

The reason comes first and the result after it, in one sentence: \textbf{\ta{வந்ததால்}}, \textbf{\ta{போனதால்}}, \textbf{\ta{இருந்ததால்}}.

\subsection*{You'll want to know: the reason after: \ta{ஏனென்றால்}, and \ta{காரணம்}}
\begin{quote}
\textbf{\ta{நான்} \ta{வரவில்லை}. \ta{ஏனென்றால்} \ta{மழை} \ta{பெய்தது}.}
\end{quote}

\begin{quote}
\emph{nāṉ varavillai. ēṉeṉṟāl maḻai peytatu.} — \textquotedblleft{}I did not come. Because it rained.\textquotedblright{}
\end{quote}

The same reason in the other order: the result first, then \textbf{\ta{ஏனென்றால்}} opens a second sentence with the cause. In speech you hear \textbf{\ta{ஏன்னா}}.

\begin{quote}
\textbf{\ta{அதற்குக்} \ta{காரணம்} \ta{மழை}.}
\end{quote}

\begin{quote}
\emph{ataṟkuk kāraṇam maḻai.} — \textquotedblleft{}The reason for that is the rain.\textquotedblright{}
\end{quote}

\textbf{\ta{அதற்கு}} (\emph{ataṟku}, \textquotedblleft{}for that\textquotedblright{}) is \textbf{\ta{அது}} with \textbf{-\ta{க்கு}}; \textbf{\ta{காரணம்}}, \textquotedblleft{}the reason\textquotedblright{}, stands in the middle, and the reason itself comes last.

\subsection*{Your turn}
Say these aloud:

\begin{itemize}
\raggedright
  \item \textquotedblleft{}Because it rained, I did not come\textquotedblright{} — \emph{maḻai peytatāl nāṉ varavillai}
  \item \textquotedblleft{}Because I had a fever, I did not go to school\textquotedblright{} — \emph{eṉakkuk kāyccal iruntatāl nāṉ paḷḷikkup pōkavillai}
  \item \textquotedblleft{}I did not come. Because it rained.\textquotedblright{} — \emph{nāṉ varavillai. ēṉeṉṟāl maḻai peytatu.}
  \item \textquotedblleft{}The reason for that is the rain\textquotedblright{} — \emph{ataṟkuk kāraṇam maḻai}
\end{itemize}

\begin{tcolorbox}[breakable,colback=teal!4,colframe=teal!35!black,title={Before you move on}]
\textquotedblleft{}Because it rained\textquotedblright{}? (\textbf{\ta{மழை} \ta{பெய்ததால்}}.) \textquotedblleft{}Because he came\textquotedblright{}? (\textbf{\ta{அவன்} \ta{வந்ததால்}}.) Which word puts the reason after the result? (\textbf{\ta{ஏனென்றால்}}.) \textquotedblleft{}The reason for that\textquotedblright{}? (\textbf{\ta{அதற்குக்} \ta{காரணம்}}.)
\end{tcolorbox}
\section[(review)]{Review — advice and reasons}

\label{lesson:TA-R287-advice-and-reasons}



\begin{quote}
Say \textquotedblleft{}I must go\textquotedblright{} and \textquotedblleft{}You must not go alone at night.\textquotedblright{} Then read \textbf{\ta{ஈரம்}} aloud. (\emph{īram}.)
\end{quote}

\subsection*{Your turn}
Say these aloud:

\begin{itemize}
\raggedright
  \item You must not eat now. Because you have a fever.
  \item Because I missed the bus, I walked.
  \item In my opinion, you must see a doctor.
  \item You don't need to come tomorrow. So you may stay at home.
  \item one piece of advice for a friend, with \ta{வேண்டும்} or \ta{கூடாது}, and the reason, with -\ta{ததால்} or \ta{ஏனென்றால்}
\end{itemize}

\begin{tcolorbox}[breakable,colback=teal!4,colframe=teal!35!black,title={Before you move on}]
\textquotedblleft{}Must\textquotedblright{}? (\textbf{to-form + \ta{வேண்டும்}}.) \textquotedblleft{}Must not\textquotedblright{}? (\textbf{to-form + \ta{கூடாது}}.) \textquotedblleft{}Because it rained\textquotedblright{}? (\textbf{\ta{மழை} \ta{பெய்ததால்}}.) \textquotedblleft{}I did not come, because …\textquotedblright{}? (\textbf{\ta{நான்} \ta{வரவில்லை}. \ta{ஏனென்றால்} …})
\end{tcolorbox}
